\documentclass[11pt]{article}
\usepackage[margin=1in]{geometry}
\usepackage{amsmath,amssymb,amsthm}
\usepackage{booktabs}
\usepackage{graphicx}
\usepackage{xcolor}
\usepackage[round,authoryear]{natbib}
\usepackage[hidelinks]{hyperref}
\usepackage{enumitem}
\setlist{itemsep=2pt,topsep=4pt}

\newcommand{\R}{\mathbb{R}}
\newcommand{\E}{\mathbb{E}}
\newcommand{\KL}{\mathrm{KL}}
\newcommand{\LN}{\mathrm{LN}}
\newcommand{\N}{\mathcal{N}}
\newcommand{\MASK}{\texttt{[MASK]}}
\newcommand{\planned}{\textsuperscript{\textcolor{gray}{\,[planned]}}}
\newtheorem{proposition}{Proposition}
\newtheorem{definition}{Definition}

\title{Generative Models of Transformer Activation Traces:\\Methods for a Trace VAE and a Prompt-Conditioned Latent Diffusion Model}
\author{SpecDec-AF (GPT-2), rev-7 / rev-8}
\date{October 2026}

\begin{document}
\maketitle

\begin{abstract}
We treat the activation trace of a frozen language model at one token position as a structured object to be modelled
generatively, in the same way that weight-space generative models treat network parameters. The trace collects every
weight--data product the model computes at that position, arranged as one fixed-schema chunk per transformer block. The
method has two stages, following latent diffusion. First, a single variational autoencoder, shared across blocks and
conditioned only on the block index, compresses each chunk to a 128-dimensional latent. Second, a diffusion transformer
over the twelve block latents of a position generates the trace conditioned on the prompt. The conditional
$p(\text{trace}\mid\text{full prefix})$ is a point mass, so we condition on a \emph{masked, left-truncated} prefix, using
the masking scheme of masked diffusion language models. This gives a family of non-degenerate conditionals, whose fully
masked member is the unconditional activation prior. It also makes ablation by marginalization available without
resampling the input. Training data is generated on the fly by the frozen teacher, with exactly reproducible resume.
We specify both stages, the data pipeline and the evaluation protocol. We report results for the VAE (\texttt{rev7B\_optd})
and the data pipeline. The latent diffusion model is specified here, but not yet trained; components marked
\planned{} have not been run.
\end{abstract}

\section{Overview and framing}
\label{sec:overview}

Our object of study is the \emph{trace} $\tau(p)$ of GPT-2 small \citep{radford2019gpt2} at token position $p$: the outputs of
every linear map and normalization in every block at that position, packed into twelve chunks of a fixed schema
(Section~\ref{sec:traces}). This follows the weight-space generative modelling line of work
\citep{soro2025d2nwg,dwf2026}. There, a network's parameters are cut into chunks with a shared schema, each chunk is
compressed by one autoencoder shared across chunks and conditioned on a chunk index, and a generative model is fitted
in the latent space. We apply the same template to activations. A trace is a weight--data product: the weights are fixed
and the data varies. The generator therefore models the distribution that the frozen network's computation induces on
its own intermediate state. Next-token prediction is a consequence of reconstructing the trace faithfully, since the last
chunk contains the input to the unembedding. It is not the training target.

The method factorizes as in latent diffusion \citep{rombach2022ldm}:
\begin{enumerate}[leftmargin=1.5em]
  \item \textbf{Trace VAE} (Section~\ref{sec:vae}): an unconditional per-block compressor
    $E: C_l \mapsto z_l \in \R^{128}$ and decoder $D:(z_l, l)\mapsto \hat C_l$, shared across the twelve blocks.
  \item \textbf{Latent generator} (Section~\ref{sec:ldm}): a diffusion transformer over $Z = (z_0,\dots,z_{11}) \in \R^{12\times128}$,
    conditioned on the prompt.
\end{enumerate}
All conditioning on context belongs to the generator. The VAE never sees the prompt, so whatever context structure the
latent carries must come from the chunk itself. A data pipeline that generates traces on the fly
(Section~\ref{sec:stream}) feeds both stages.

\section{Traces}
\label{sec:traces}

\paragraph{Teacher.} GPT-2 small (12 blocks, $d_{\text{model}} = 768$), frozen, in evaluation mode, fp32. For a token
sequence $x_{1:L}$, the residual stream at position $p$ is $r_0(p) = W_E x_p + W_P p$ and
\begin{equation}
  a_l = \mathrm{Attn}_l\big(\LN_{1,l}(r_l)\big),\qquad
  m_l = \mathrm{MLP}_l\big(\LN_{2,l}(r_l + a_l)\big),\qquad
  r_{l+1} = r_l + a_l + m_l,
\end{equation}
for $l = 0,\dots,11$. The logits are $W_U\,\LN_f(r_{12})$. Attention is causal, so $\tau(p)$ is a deterministic function of
$x_{1:p}$.

\paragraph{Chunk schema.} The chunk of block $l$ at position $p$ is $C_l(p) \in \R^{9984}$, the concatenation of eight
slots (Table~\ref{tab:schema}). The two boundary slots carry the embedding at $l = 0$ and the final-norm output at
$l = 11$, and are zero for other blocks. Every block therefore has the same schema, and the block index is the only thing
that distinguishes one chunk from another. The unembedding input is slot~7 of $C_{11}$, so the decoded trace yields
next-token logits directly.

\begin{table}[h]
  \centering\small
  \begin{tabular}{@{}rlrl@{}}
    \toprule
    Slot & Module output & Width & Content \\
    \midrule
    0 & \texttt{transformer.drop} & 768 & $r_0$ at $l=0$; zeros otherwise \\
    1 & \texttt{h[l].ln\_1} & 768 & $\LN_{1,l}(r_l)$ \\
    2 & \texttt{h[l].attn.c\_attn} & 2304 & $Q, K, V$ \\
    3 & \texttt{h[l].attn.c\_proj} & 768 & $a_l$ \\
    4 & \texttt{h[l].ln\_2} & 768 & $\LN_{2,l}(r_l + a_l)$ \\
    5 & \texttt{h[l].mlp.c\_fc} & 3072 & MLP pre-activation \\
    6 & \texttt{h[l].mlp.c\_proj} & 768 & $m_l$ \\
    7 & \texttt{transformer.ln\_f} & 768 & $\LN_f(r_{12})$ at $l=11$; zeros otherwise \\
    \midrule
    & & 9984 & \\
    \bottomrule
  \end{tabular}
  \caption{Chunk schema (one chunk per block, twelve per position).}
  \label{tab:schema}
\end{table}

\paragraph{Internal consistency.} A trace is heavily redundant under the frozen weights. With HF \texttt{Conv1D}
convention $y = xW + b$, slots 2, 5 and 6 are affine (or GELU-then-affine) functions of slots 1, 4 and 5 of the same block:
\begin{equation}
  \text{c\_attn}_l = \text{ln\_1}_l W^{\text{attn}}_l + b^{\text{attn}}_l,\quad
  \text{c\_fc}_l = \text{ln\_2}_l W^{\text{fc}}_l + b^{\text{fc}}_l,\quad
  m_l = \mathrm{gelu}(\text{c\_fc}_l) W^{\text{proj}}_l + b^{\text{proj}}_l .
  \label{eq:intra}
\end{equation}
Rebuilding the residual stream from slots 0, 3 and 6 gives cross-block identities:
$\text{ln\_1}_l = \LN_{1,l}(r_l)$, $\text{ln\_2}_l = \LN_{2,l}(r_l + a_l)$ and $\text{ln\_f} = \LN_f(r_{12})$.
The attention output $a_l$ is the one slot that cannot be checked from one position, because it mixes in other positions'
keys and values. We use the relative violation $\lVert \text{lhs} - \text{rhs}\rVert / \lVert\text{rhs}\rVert$ as a
reference-free measure of how close a generated trace is to a trace the teacher could have computed. On real traces the
violation sits at the fp16 storage floor ($\approx 2.6\times10^{-4}$ cross-block).

\paragraph{Positions and corpus.} Sequences are non-overlapping packed 256-token windows of WikiText-103 train
\citep{merity2016wikitext}, tokenized without special tokens, with document boundaries ignored. From each sequence we draw
16 positions uniformly without replacement from $[32, 255]$, so every target has at least 32 tokens of context. Draws are
a pure function of (seed, global sequence index, stream), so any position can be regenerated. Val is a fixed set of 360
sequences (global indices 5890--6249); every val metric below uses the same first 5{,}000 val positions.

\section{Trace VAE}
\label{sec:vae}

\subsection{Model}

\paragraph{Normalization.} Chunk elements vary in scale by orders of magnitude, across slots and across depth. We fit a
fixed per-block, per-element affine normalization $N_l(C) = (C - \bar C_l)/\sigma_l$ once, with a streaming
Welford pass over 16k calibration positions. These are drawn from the training position distribution with a separate
random stream. $\sigma_l$ is clamped at $10^{-6}$, and a mask $M_l$ removes the zero-padded boundary slots from every loss
term. The statistics are frozen and stored in every checkpoint.

\paragraph{Encoder and decoder.} One encoder and one decoder are shared across the twelve blocks. The encoder is
block-agnostic and unconditioned. It maps $N_l(C)$ through three Linear--LayerNorm--GELU layers ($9984 \to 2048 \to 1024
\to 512$) to a diagonal Gaussian posterior $q(z\mid C) = \N(\mu_E, \mathrm{diag}\,e^{\lambda_E})$, $z\in\R^{128}$. The
decoder is conditioned only on the block index: it maps $[z;\, e_l]$, with a learned 64-dimensional block embedding $e_l$,
through the mirrored tower ($192 \to 512 \to 1024 \to 2048$) and a linear head to $\R^{9984}$. The model has no context,
position or prompt input. A trace's latent is the stack of its twelve block latents, $Z \in \R^{12\times128}$ (1{,}536
numbers per position).

\paragraph{Output space.} The decoder emits \emph{raw} activations. The loss is a $\sigma$-weighted squared error in raw
space, which equals the MSE in normalized space up to the mask:
\begin{equation}
  \mathcal{L}_{\text{rec}} = \frac{\sum_j M_{l,j}\,\sigma_{l,j}^{-2}\,\big(\hat C_{l,j} - C_{l,j}\big)^2}{\sum_j M_{l,j}} .
\end{equation}
Decoded traces can then be fed to the teacher or the unembedding with no inverse transform. We compared this against
decoding in normalized space with a raw-space MSE, under a pre-registered rule (teacher top-1 after 60k steps). The raw
output won, 0.483 vs 0.411. The final linear layer is initialized with each output row scaled by the
block-averaged $\sigma$, so that outputs start at the target scale.

\paragraph{Objective.} We use the standard ELBO with a linearly annealed KL weight \citep{kingma2014vae,bowman2016generating}
and free bits \citep{kingma2016iaf}:
\begin{equation}
  \mathcal{L} = \mathcal{L}_{\text{rec}}\big(D(z, l), C\big) + \beta_s \cdot \frac{1}{128}\sum_{k=1}^{128}
  \max\!\Big(\phi,\ \E_{\text{batch}}\,\KL\big(q(z_k\mid C)\,\Vert\,\N(0,1)\big)\Big),
  \qquad z \sim q(z\mid C),
\end{equation}
with $\beta_s = \beta_{\max}\min(1, s/S_\beta)$, $\beta_{\max} = 0.01$ and $\phi = 0.1$ nats. Each training item is one
(position, block) pair, and a batch mixes blocks.

\begin{table}[h]
  \centering\small
  \begin{tabular}{@{}ll@{}}
    \toprule
    Parameter & Value (\texttt{rev7B\_optd}) \\
    \midrule
    Data & 500{,}768 positions (31.3k sequences $\times$ 16), i.e.\ 6.0M (position, block) items \\
    Optimizer & Adam, lr $5\times10^{-4}$, 500-step linear warmup, grad-norm clip 1.0 \\
    Batch / steps & 256 items, 200k steps (about 8.5 epochs) \\
    KL & $\beta_{\max} = 0.01$, linear ramp over $S_\beta = 80$k steps, free bits 0.1 nats/dim \\
    Latent & $d_{\text{latent}} = 128$ per block ($12\times128$ per position) \\
    \bottomrule
  \end{tabular}
  \caption{VAE training configuration.}
  \label{tab:vae-config}
\end{table}

\subsection{Evaluation}
\label{sec:vae-eval}

All metrics use the 5{,}000 fixed val positions. Latent conditions:
\begin{itemize}[leftmargin=1.5em]
  \item \texttt{qz\_mean} and \texttt{qz\_sample}: decode $\mu_E$ or a posterior sample;
  \item \texttt{prior}: decode $z\sim\N(0,I)$;
  \item \texttt{wrong\_z}: decode a different position's latent.
\end{itemize}
The last two are chance-level controls.

\begin{itemize}[leftmargin=1.5em]
  \item \textbf{Downstream fidelity.} Let $\hat P = \mathrm{softmax}(W_U \hat C_{11}^{(7)})$ be the decoded terminal slot read
    out by the unembedding, and $P$ the teacher distribution from the real trace. We report top-1 agreement, top-5 overlap and
    $\KL(P\,\Vert\,\hat P)$.
  \item \textbf{Reconstruction.} Per-(block, slot) cosine and normalized MSE.
  \item \textbf{Consistency.} Median intra- and cross-block violations of Eq.~\eqref{eq:intra} and the cross-block identities,
    against the fp16 floor of real traces.
  \item \textbf{Latent usage.} Active units (posterior-mean variance $> 10^{-2}$), KL per block, and the participation ratio of the
    aggregate posterior.
  \item \textbf{In-model substitution.} A block's reconstruction is put back into the teacher's forward pass at a single
    position, and the divergence is measured at every downstream block and at the logits. This is specified separately in
    \texttt{docs/block\_substitution\_method.tex}.
\end{itemize}

\subsection{Results}

\begin{table}[h]
  \centering\small
  \begin{tabular}{@{}lr@{}}
    \toprule
    Metric (val, 5{,}000 positions) & \texttt{rev7B\_optd} \\
    \midrule
    Teacher top-1, \texttt{qz\_mean} (train) & 0.518 (0.525) \\
    Teacher top-1, \texttt{qz\_sample} & 0.517 \\
    Top-5 overlap & 0.566 \\
    $\KL(P\,\Vert\,\hat P)$ (nats) & 0.89 \\
    Top-1, \texttt{prior} / \texttt{wrong\_z} & 0.032 / 0.044 \\
    Recon cosine, mean / worst block & 0.938 / 0.880 \\
    Cross-block consistency, median (fp16 floor $2.6\times10^{-4}$) & 0.350 \\
    Intra-block consistency, median (fp16 floor $2.8\times10^{-4}$) & 0.165 \\
    Active units & 1{,}536 / 1{,}536 \\
    KL, summed over blocks (nats) & 2{,}810 \\
    \bottomrule
  \end{tabular}
  \caption{Trace VAE results. The train/val top-1 gap is 0.6\,pp.}
  \label{tab:vae-results}
\end{table}

These results bear directly on the second stage:
\begin{itemize}[leftmargin=1.5em]
  \item \textbf{No low-dimensional shortcut.} All 128 dimensions are active in all twelve blocks. The participation ratio of
    the aggregate posterior is 72--100 per block, and 99--106 dimensions are needed for 90\% of the variance.
  \item \textbf{The aggregate posterior is far from the prior.} Decoding $\N(0,I)$ gives chance-level top-1, so the
    distribution has to be learned by a generator; this is the regime latent diffusion is designed for.
  \item \textbf{Weak linear coupling across depth.} The mean absolute correlation of $\mu$ between blocks is 0.052, so the
    joint structure of the twelve latents is largely non-Gaussian. This motivates a joint model over depth rather than
    independent per-block models.
  \item \textbf{The VAE sets the ceiling.} \texttt{qz\_mean} defines the \emph{VAE floor} that the generator's decoded
    samples can at best reach. Generator metrics are reported as a gap to this floor.
\end{itemize}

\section{Streaming data pipeline}
\label{sec:stream}

A static activation cache costs about 240\,MB per thousand positions (112\,GiB for 500k), so the 10--100M positions we
expect the generator to need are out of reach on shared scratch storage. Instead, the frozen teacher runs on the training
GPU, and training batches are produced on the fly. The only stored artefacts are a pre-tokenized corpus (118M tokens, 470\,MB),
the normalization statistics, and the fixed val shards.

\paragraph{Plan.} Sequence $i$ of the corpus is the slice $[256i, 256(i+1))$ of the pre-tokenized array. This is exactly
the packing used to build the cache, so sequence ids agree across the two paths, and any sequence can be read directly.
Let $\mathcal{I}$ be the train ids, which is every sequence outside the val range (460{,}265 of 460{,}625).
\begin{itemize}[leftmargin=1.5em]
  \item Pass $k$ is the permutation $\pi_k = \mathrm{perm}(\text{seed}, k)$ of $\mathcal{I}$, cut into chunks of $K = 256$
    sequences. Global chunk $c$ is chunk $c \bmod \lfloor|\mathcal{I}|/K\rfloor$ of pass $k = \lfloor c / \lfloor|\mathcal{I}|/K\rfloor\rfloor$.
  \item Positions are re-drawn in every pass with random stream $2 + k$, so a pass yields about 7.4M new (sequence, position)
    pairs.
  \item The teacher runs on chunk $c$ in micro-batches, and its traces are rounded to fp16. This matches the distribution of
    the cache the VAE was trained on.
  \item The resulting $4{,}096$ positions are permuted with $\mathrm{perm}(\text{seed}, c)$ and served in fixed-size batches.
    For the VAE, items are $(\text{position}, \text{block})$ pairs, so one chunk gives 192 batches of 256.
  \item The latent-generator consumer encodes four chunks with the frozen VAE and shuffles the 16{,}384 resulting
    positions. Each item is $(\mu_E, \lambda_E)$ for all twelve blocks, plus the source sequence, the position and the
    sequence id.
\end{itemize}

\begin{proposition}[Exact resume]
Batch $j$ of buffer $c$ is a deterministic function of (corpus, seed, $c$, $j$) and the frozen teacher (and frozen VAE).
Checkpointing the cursor $(c, j)$, alongside the optimizer and RNG state, is therefore enough to resume training exactly, with no buffer state and no data-loader
state.
\end{proposition}
\noindent The only caveat is GPU kernel nondeterminism in the teacher forward, which changes about 0.8\% of trace elements
by one fp16 rounding step. On CPU, resume is bit-exact, and this is tested: 10 steps, resume, then 10 more steps gives
the same weights as 20 uninterrupted steps, across buffer boundaries.

\paragraph{Validation.} On one V100, against cache shards 0000 (Stage A), 0100 (val) and 0504 (Stage B):
\begin{itemize}[leftmargin=1.5em]
  \item Tokens are identical.
  \item Live traces are 99.2\% bit-identical to the cache, and otherwise within one fp16 step (relative error $2\times10^{-5}$).
  \item Encoder means agree to $3\times10^{-4}$.
\end{itemize}
A 5k-step VAE run reached val reconstruction loss 0.5816 on the stream and 0.5841 on the cache, with the two curves
agreeing at every 1k-step checkpoint. The stream run took 113\,s against 482\,s for the disk-bound cache run. The
teacher produces about 1{,}760 positions/s, and the latent consumer about 2{,}100 positions/s.

\section{Prompt-conditioned latent diffusion}
\label{sec:ldm}

\subsection{Target distribution}

\paragraph{Degeneracy.} Since $\tau(T)$ is a deterministic function of $x_{1:T}$, the natural conditional
$p(Z_T\mid x_{1:T})$ is a point mass. Fitted with a diffusion model, it reduces to regression with extra steps: the
diversity, the prior and the score geometry that make diffusion models useful objects of study all vanish.

\paragraph{Masked, left-truncated conditioning.} We instead condition on a partially observed window. Fix a window
length $W = 128$ ending at $T$. A visibility mask $S \subseteq \{T-W+1, \dots, T\}$ selects which window tokens are shown;
all other window slots, and all tokens before the window, are hidden. The target is
\begin{equation}
  p(Z_T \mid x_S, T) = \int \delta\big(Z_T - E(\tau(x_{1:T}))\big)\; p\big(x_{\bar S}\mid x_S, T\big)\, dx_{\bar S},
  \label{eq:target}
\end{equation}
where $\bar S$ is every unobserved position $\le T$ and $E$ applies the encoder blockwise. This is the push-forward of
the corpus's conditional distribution over the unseen context, through the teacher and the encoder. It is non-degenerate
whenever $\bar S$ is nonempty. Masks are drawn independently of the data. A conditional model trained on pairs
$(Z_T, x_S)$, where $x_S$ is the true sequence with the mask applied, therefore targets Eq.~\eqref{eq:target} exactly.
Hidden slots are filled with the \MASK{} token of a masked diffusion language model \citep{sahoo2024mdlm}, which is
pretrained to read partially observed text. Masking is the mechanism; the encoder supplies the vocabulary for it.

\paragraph{Mask distribution.} For each training example:
\begin{itemize}[leftmargin=1.5em]
  \item \textbf{null} ($p = 0.10$): every slot is masked, including $x_T$. This gives the unconditional activation prior
    $p(Z)$ and the null branch for classifier-free guidance \citep{ho2022cfg}.
  \item \textbf{suffix} ($p = 0.75$): the last $m+1$ tokens $x_{T-m:T}$ are visible. The context length $m$ is drawn from
    $\{0\}$ w.p.\ 0.15, $\mathcal{U}\{1..8\}$ w.p.\ 0.45, $\mathcal{U}\{9..64\}$ w.p.\ 0.25 and $W-1$ w.p.\ 0.15.
  \item \textbf{holes} ($p = 0.15$): every window slot except $x_T$ is masked independently at rate $r\sim\mathcal{U}(0, 0.5)$.
\end{itemize}
A single model thus spans a family of conditionals: the unconditional prior ($S = \varnothing$), the current token alone
($m = 0$), and conditionals that approach full-context regression ($m = W-1$). No token after $T$ ever enters the window.

\paragraph{Attribution by marginalization.} Eq.~\eqref{eq:target} gives token-level attribution a direct
interventional reading. For $i \in S$, comparing $p(Z_T\mid x_S)$ with $p(Z_T\mid x_{S\setminus\{i\}})$ is a
\emph{resample ablation}: $x_i$ is replaced by draws from its corpus conditional given the remaining visible tokens, and
the induced trace distribution is compared with the original. Both distributions are sampled from the same generator by
changing only the mask. There is no substitute text to sample, and no off-distribution zero or mean ablation.

\subsection{Architecture}

The denoiser follows the cross-attention diffusion transformer used in text-to-image models
\citep{peebles2023dit,chen2024pixart}. These are the systems on which most mechanistic studies of diffusion models have
been done (sparse autoencoders, cross-attention attribution, parameter localization). We keep our design close to them so
that those tools transfer.
\begin{itemize}[leftmargin=1.5em]
  \item \textbf{Tokens.} One token per block: $z_{t,l} \in \R^{128}$ is projected to width $d$, and a learned depth embedding
    is added. Self-attention runs over the 12 depth tokens of one position.
  \item \textbf{Prompt encoder.} The token embedding and \emph{first} DDiT block of the pretrained MDLM-OWT model
    \citep{sahoo2024mdlm}, frozen. This model uses the GPT-2 tokenizer plus a \MASK{} token, and is trained on OpenWebText.
    The encoder is deliberately shallow:
    \begin{itemize}[leftmargin=1.2em]
      \item It contextualizes the visible tokens and understands \MASK{}.
      \item It is far too shallow to reproduce the teacher's computation, so it cannot hand the generator the trace.
      \item MDLM uses rotary position embeddings only, so the absolute position $T$ is supplied separately.
      \item Its blocks are modulated by a time embedding; with time conditioning off this is the constant
        $\mathrm{silu}(\text{sigma\_map}(0))$.
    \end{itemize}
  \item \textbf{Block.} Each layer applies, in order:
    \begin{enumerate}[leftmargin=1.2em]
      \item adaLN-Zero self-attention;
      \item unmodulated cross-attention, with queries from the depth tokens and keys/values from a linear projection of the
        encoder states;
      \item an adaLN-Zero SwiGLU MLP \citep{shazeer2020glu}.
    \end{enumerate}
    The modulation vector is per token,
    \[ c_{\text{tok}} = \text{emb}(t_{\text{tok}}) + \text{emb}(T) + \text{proj}(\text{mean-pooled encoder states}). \]
    The final layer is zero-initialized, so the model outputs $v \equiv 0$ at initialization.
  \item \textbf{No joint attention.} Joint text--latent attention, as in MM-DiT \citep{esser2024sd3}, ranks higher in image
    generation. We exclude it on purpose. It would let the representation of the prompt be rewritten by the activations we
    want to interpret, so that any attribution from prompt to trace would partly measure the generator's own mixing.
    Keeping the encoder frozen and the conditioning one-way (cross-attention) keeps the prompt side fixed.
  \item \textbf{Size.} DiT-S to start: depth 12, $d = 384$, 6 heads.
\end{itemize}

\subsection{Objective and sampling}

We train with rectified flow \citep{liu2023rectified,lipman2023flow}:
\begin{equation}
  z_t = (1-t)\,z_0 + t\,\varepsilon,\quad \varepsilon\sim\N(0, I),\qquad
  \mathcal{L} = \E\,\big\lVert v_\theta(z_t, t, x_S, T) - (\varepsilon - z_0)\big\rVert^2,
\end{equation}
with $t = \mathrm{sigmoid}(u)$, $u\sim\N(0,1)$ (logit-normal; \citealp{esser2024sd3}). The clean-latent estimate
$\hat z_0 = z_t - t\,v_\theta$ is the single accessor all analysis code uses. With probability 0.1 per batch, each of the
12 depth tokens gets an independent $t$, following Diffusion Forcing \citep{chen2024diffusionforcing}. This keeps
depth-wise partial-noise queries in distribution at almost no cost. Sampling is an Euler ODE from $t=1$ to 0 (50 steps by
default), with classifier-free guidance $v = v_{\text{null}} + w(v_{\text{cond}} - v_{\text{null}})$. We keep an EMA
(0.9999) and evaluate on it.

\paragraph{Latent coordinates.} The generator works directly in the encoder's coordinates. There is no per-dimension
standardization and no global scale factor. No transform is coordinate-free, and the VAE's KL term already fixes the
scale of its own coordinates near the prior. Any whitening would instead encode the statistics of one particular
training sample. Jacobian-based analyses therefore include the encoder's affine terms; we accept that.

\paragraph{Training targets.} The target $z_0$ can be the encoder mean $\mu_E$ or a posterior sample. The posterior is
narrow ($\sigma \approx 0.16$ against a per-dimension aggregate standard deviation of about 0.9), so the two differ mainly
in augmentation. Streaming removes most of the need for augmentation. The choice is settled by a pre-registered
head-to-head on the stream: two 50k-step runs, scored on decoded teacher KL averaged over $m$. A run that diverges or fails
the memorization screen is disqualified, and ties within 3\% go to $\mu_E$.\planned

\subsection{Evaluation protocol\planned}

All conditions are evaluated on the 5{,}000 VAE val positions with fixed seeds: null; suffix with
$m\in\{0, 2, 8, 32, W-1\}$; holes with $r = 0.1$; and shuffled (another position's full window).
\begin{itemize}[leftmargin=1.5em]
  \item \textbf{Per-sample fidelity.} We decode each sample and score it against the real trace with the VAE metrics of
    Section~\ref{sec:vae-eval}. The headline number is the \emph{gap to the VAE floor} $D(\mu_E(\tau))$, reported by $m$. It
    should shrink as $m$ grows, and stay large under \texttt{shuffled}.
  \item \textbf{Consistency} of decoded samples, against both the VAE floor and the fp16 floor.
  \item \textbf{Distributional fit.} We compute Fréchet distance \citep{heusel2017fid} and kernel inception distance
    \citep{binkowski2018kid} between samples and real latents, pooled over conditions, in three spaces:
    \begin{itemize}[leftmargin=1.2em]
      \item per-block latents (128-d);
      \item the joint latent (1{,}536-d);
      \item the decoded residual (768-d).
    \end{itemize}
    Every condition has the same marginal target, the val distribution, so the pooled statistic is a calibration test even
    though each condition has only one sample per position. KID is primary: FD is biased at $N = 5$k in 1{,}536 dimensions.
    Both are bracketed by a data floor (val against train) and a prior ceiling ($\N(0,I)$ against val).
  \item \textbf{Diversity against $m$.} Pairwise distance and decoded top-1 agreement among $K = 4$ samples per condition.
    These should fall monotonically as $m$ increases.
  \item \textbf{Guidance sweep.} CFG with $w \in\{1, 1.5, 2, 3\}$.
  \item \textbf{Memorization screen.} For each null sample, the distance to its nearest neighbour in a 100k training
    subsample, compared with the same statistic for val latents.
  \item \textbf{Generalization by disjoint halves} \citep{kadkhodaie2024generalization}. Two models are trained identically on
    the even and odd sequence ids. Both are sampled from shared noise under identical conditions, and we report the paired
    distance relative to the distance between independent samples:
    \begin{itemize}[leftmargin=1.2em]
      \item if the paired samples nearly coincide, the models have learned the density;
      \item if they diverge, each model has memorized its half.
    \end{itemize}
    This test is what licenses later geometric analyses of the learned score.
\end{itemize}

\paragraph{Baselines.\planned}
\begin{itemize}[leftmargin=1.5em]
  \item An unconditional deep SwiGLU MLP over the flattened 1{,}536-d latent, in the style of generative activation priors
    \citep{luo2026glp}, evaluated under null only.
  \item The main architecture with diagonal self-attention. The twelve blocks are then independent given the condition,
    which isolates what joint modelling across depth buys, particularly for cross-block consistency.
\end{itemize}

\section{Status and limitations}

\begin{itemize}[leftmargin=1.5em]
  \item \textbf{Built and evaluated:}
    \begin{itemize}[leftmargin=1.2em]
      \item the chunk schema and the cache;
      \item the trace VAE (Section~\ref{sec:vae});
      \item block substitution;
      \item the streaming pipeline with both consumers (Section~\ref{sec:stream}).
    \end{itemize}
  \item \textbf{Specified, not yet run:} the whole of Section~\ref{sec:ldm}, including the MDLM block-0 encoder, which needs
    a functional check against the full MDLM before anything is built on it.
  \item \textbf{The VAE floor bounds every generator result.} About half of val positions decode to a different argmax than
    the teacher's (top-1 0.518). Decoded traces violate the teacher's cross-block identities about 1{,}350$\times$ above the
    fp16 floor. A larger latent ($d = 256$) is the next lever. It will be chosen with a rule that accounts for how
    diffusable the latent is, not only for reconstruction \citep{yao2025vavae,skorokhodov2025diffusability}.
  \item \textbf{Single-position traces.} A trace omits the keys and values of other positions, which is why $a_l$ cannot be
    checked for consistency. Multi-position generation (grids over tokens $\times$ depth) is future work.
  \item \textbf{One corpus, one teacher.} WikiText-103 and GPT-2 small. The stream makes changing the corpus cheap; scaling
    the teacher is not.
\end{itemize}

\begin{thebibliography}{99}\small
\setlength{\itemsep}{1pt}

\bibitem[Bi{\'n}kowski et~al.(2018)]{binkowski2018kid}
M.~Bi{\'n}kowski, D.~J. Sutherland, M.~Arbel, A.~Gretton.
\newblock Demystifying {MMD} {GAN}s.
\newblock \emph{ICLR}, 2018. arXiv:1801.01401.

\bibitem[Bowman et~al.(2016)]{bowman2016generating}
S.~R. Bowman, L.~Vilnis, O.~Vinyals, A.~M. Dai, R.~Jozefowicz, S.~Bengio.
\newblock Generating sentences from a continuous space.
\newblock \emph{CoNLL}, 2016. arXiv:1511.06349.

\bibitem[Chen et~al.(2024a)]{chen2024diffusionforcing}
B.~Chen, D.~Mart{\'i} Mons{\'o}, Y.~Du, M.~Simchowitz, R.~Tedrake, V.~Sitzmann.
\newblock Diffusion forcing: Next-token prediction meets full-sequence diffusion.
\newblock \emph{NeurIPS}, 2024. arXiv:2407.01392.

\bibitem[Chen et~al.(2024b)]{chen2024pixart}
J.~Chen et~al.
\newblock {PixArt}-$\alpha$: Fast training of diffusion transformer for photorealistic text-to-image synthesis.
\newblock \emph{ICLR}, 2024. arXiv:2310.00426.

\bibitem[DWF(2026)]{dwf2026}
S.~Biggs et~al.
\newblock {DWF}: chunked-{VAE} latent flow matching over network weights.
\newblock arXiv:2601.05052, 2026. \textcolor{red}{[title/authors to confirm]}

\bibitem[Esser et~al.(2024)]{esser2024sd3}
P.~Esser et~al.
\newblock Scaling rectified flow transformers for high-resolution image synthesis.
\newblock \emph{ICML}, 2024. arXiv:2403.03206.

\bibitem[Heusel et~al.(2017)]{heusel2017fid}
M.~Heusel, H.~Ramsauer, T.~Unterthiner, B.~Nessler, S.~Hochreiter.
\newblock {GAN}s trained by a two time-scale update rule converge to a local {N}ash equilibrium.
\newblock \emph{NeurIPS}, 2017. arXiv:1706.08500.

\bibitem[Ho \& Salimans(2022)]{ho2022cfg}
J.~Ho, T.~Salimans.
\newblock Classifier-free diffusion guidance.
\newblock arXiv:2207.12598, 2022.

\bibitem[Kadkhodaie et~al.(2024)]{kadkhodaie2024generalization}
Z.~Kadkhodaie, F.~Guth, E.~P. Simoncelli, S.~Mallat.
\newblock Generalization in diffusion models arises from geometry-adaptive harmonic representations.
\newblock \emph{ICLR}, 2024. arXiv:2310.02557.

\bibitem[Kingma \& Welling(2014)]{kingma2014vae}
D.~P. Kingma, M.~Welling.
\newblock Auto-encoding variational {B}ayes.
\newblock \emph{ICLR}, 2014. arXiv:1312.6114.

\bibitem[Kingma et~al.(2016)]{kingma2016iaf}
D.~P. Kingma, T.~Salimans, R.~Jozefowicz, X.~Chen, I.~Sutskever, M.~Welling.
\newblock Improved variational inference with inverse autoregressive flow.
\newblock \emph{NeurIPS}, 2016. arXiv:1606.04934.

\bibitem[Lipman et~al.(2023)]{lipman2023flow}
Y.~Lipman, R.~T.~Q. Chen, H.~Ben-Hamu, M.~Nickel, M.~Le.
\newblock Flow matching for generative modeling.
\newblock \emph{ICLR}, 2023. arXiv:2210.02747.

\bibitem[Liu et~al.(2023)]{liu2023rectified}
X.~Liu, C.~Gong, Q.~Liu.
\newblock Flow straight and fast: Learning to generate and transfer data with rectified flow.
\newblock \emph{ICLR}, 2023. arXiv:2209.03003.

\bibitem[Luo et~al.(2026)]{luo2026glp}
G.~Luo, J.~Feng, T.~Darrell, A.~Radford, J.~Steinhardt.
\newblock Learning a generative meta-model of {LLM} activations.
\newblock \emph{ICML}, 2026. arXiv:2602.06964.

\bibitem[Merity et~al.(2017)]{merity2016wikitext}
S.~Merity, C.~Xiong, J.~Bradbury, R.~Socher.
\newblock Pointer sentinel mixture models.
\newblock \emph{ICLR}, 2017. arXiv:1609.07843.

\bibitem[Peebles \& Xie(2023)]{peebles2023dit}
W.~Peebles, S.~Xie.
\newblock Scalable diffusion models with transformers.
\newblock \emph{ICCV}, 2023. arXiv:2212.09748.

\bibitem[Radford et~al.(2019)]{radford2019gpt2}
A.~Radford, J.~Wu, R.~Child, D.~Luan, D.~Amodei, I.~Sutskever.
\newblock Language models are unsupervised multitask learners.
\newblock OpenAI technical report, 2019.

\bibitem[Rombach et~al.(2022)]{rombach2022ldm}
R.~Rombach, A.~Blattmann, D.~Lorenz, P.~Esser, B.~Ommer.
\newblock High-resolution image synthesis with latent diffusion models.
\newblock \emph{CVPR}, 2022. arXiv:2112.10752.

\bibitem[Sahoo et~al.(2024)]{sahoo2024mdlm}
S.~S. Sahoo et~al.
\newblock Simple and effective masked diffusion language models.
\newblock \emph{NeurIPS}, 2024. arXiv:2406.07524.

\bibitem[Shazeer(2020)]{shazeer2020glu}
N.~Shazeer.
\newblock {GLU} variants improve transformer.
\newblock arXiv:2002.05202, 2020.

\bibitem[Skorokhodov et~al.(2025)]{skorokhodov2025diffusability}
I.~Skorokhodov et~al.
\newblock Improving the diffusability of autoencoders.
\newblock \emph{ICML}, 2025. arXiv:2502.14831.

\bibitem[Soro et~al.(2025)]{soro2025d2nwg}
B.~Soro, B.~Andreis, H.~Lee, W.~Jeong, S.~Chong, F.~Hutter, S.~J. Hwang.
\newblock Diffusion-based neural network weights generation.
\newblock \emph{ICLR}, 2025. arXiv:2402.18153.

\bibitem[Yao et~al.(2025)]{yao2025vavae}
J.~Yao, B.~Yang, X.~Wang.
\newblock Reconstruction vs.\ generation: Taming optimization dilemma with latent diffusion models.
\newblock \emph{CVPR}, 2025. arXiv:2501.01423.

\end{thebibliography}

\end{document}
